% Idioma y codificación
%\usepackage[spanish, es-tabla, es-notilde]{babel}       %es-tabla para que se titule "Tabla"
\usepackage[utf8]{inputenc}

% Márgenes
\usepackage[a4paper,top=3cm,bottom=2.5cm,left=3cm,right=3cm]{geometry}

% Comentarios de bloque
\usepackage{verbatim}

% Paquetes de links
\usepackage{url}                    % redirecciona a la web

% Más opciones para enumeraciones
\usepackage{enumitem}

% Personalizar la portada
\usepackage{titling}

% Paquetes de tablas
\usepackage{multirow}

% Para añadir el símbolo de euro
\usepackage{eurosym}

%------------------------------------------------------------------------

%Paquetes de figuras
\usepackage{caption}
\usepackage{subcaption} % Figuras al lado de otras
\usepackage{float}      % Poner figuras en el sitio indicado H.


% Paquetes de imágenes
\usepackage{graphicx}       % Paquete para añadir imágenes
\usepackage{transparent}    % Para manejar la opacidad de las figuras

% Paquete para usar colores
\usepackage[dvipsnames, table, xcdraw]{xcolor}
\usepackage{pagecolor}      % Para cambiar el color de la página

% Habilita tamaños de fuente mayores
\usepackage{fix-cm}

% Para los gráficos
\usepackage{tikz}
\usepackage{forest}

% Para poder situar los nodos en los grafos
\usetikzlibrary{positioning}

% Para ajustar las tablas
\usepackage{adjustbox}

% Para recuadrar solucions
\usepackage[most]{tcolorbox}
%\usepackage[breakable]{tcolorbox}
\tcbuselibrary{breakable, skins}

%------------------------------------------------------------------------

% Paquetes de matemáticas
\usepackage{mathtools, amsfonts, amssymb, mathrsfs}
\usepackage[makeroom]{cancel}     % Simplificar tachando
\usepackage{polynom}    % Divisiones y Ruffini
%\usepackage{units} % Para poner fracciones diagonales con \nicefrac

\usepackage{pgfplots}   %Representar funciones
\pgfplotsset{compat=1.18}  % Versión 1.18

\usepackage{tikz-cd}    % Para usar diagramas de composiciones
\usetikzlibrary{calc}   % Para usar cálculo de coordenadas en tikz

%Definición de teoremas, etc.
\usepackage{amsthm}

\usepackage{nicematrix}

\usepackage{longtable}

\usepackage{booktabs}
\usepackage{siunitx}
\usepackage{nicefrac}
%\swapnumbers   % Intercambia la posición del texto y de la numeración
\usepackage{xparse} % needed for optional argument parsing
\usepackage[hidelinks]{hyperref}    % Permite enlaces

% Save the original proof environment
\let\oldproof\proof
\let\endoldproof\endproof

% Redefine proof to accept optional parentheses
\RenewDocumentEnvironment{proof}{O{}} % optional argument
{
  \if\relax\detokenize{#1}\relax
    \oldproof % no argument, just "Proof."
  \else
    \oldproof[Proof (#1)] % argument given, "Proof (Sketch)."
  \fi
}
{
  \endoldproof
}

\theoremstyle{plain}

\makeatletter
\@ifclassloaded{article}{
  \newtheorem{teo}{Teorema}[section]
}{
  \newtheorem{teo}{Teorema}[chapter]  % Se resetea en cada chapter
}
\makeatother

\makeatletter
\@ifclassloaded{article}{
  \newtheorem{theo}{Theorem}[section]
}{
  \newtheorem{theo}{Theorem}[chapter]  % Se resetea en cada chapter
}
\makeatother
\makeatletter
\@ifclassloaded{article}{
  \newtheorem{lemma}{Lemma}[section]
}{
  \newtheorem{lemma}{Lemma}[chapter]  % Se resetea en cada chapter
}
\makeatother
\makeatletter
\@ifclassloaded{article}{
  \newtheorem{corollary}{Corollary}[section]
}{
  \newtheorem{corollary}{Corollary}[chapter]  % Se resetea en cada chapter
}
\makeatother
\makeatletter
\@ifclassloaded{article}{
  \newtheorem{pro}{Proposition}[section]
}{
  \newtheorem{pro}{Proposition}[chapter]  % Se resetea en cada chapter
}
\makeatother
\providecommand*{\corollaryautorefname}{Corollary}
\providecommand*{\theoautorefname}{Theorem}
\providecommand*{\lemmaautorefname}{Lemma}
\providecommand*{\proautorefname}{Proposition}

\newtheorem{coro}{Corolario}[teo]           % Se resetea en cada teorema
\newtheorem{prop}[teo]{Proposición}         % Usa el mismo contador que teorema
\newtheorem{lema}[teo]{Lema}                % Usa el mismo contador que teorema
%\newtheorem{lemma}[lemma]{Lemma}                % Usa el mismo contador que teorema
\newtheorem*{lema*}{Lema}
\newtheorem*{notation}{Notation}
\newtheorem*{suggestion}{}

\theoremstyle{remark}
\newtheorem*{observacion}{Observación}
\newtheorem*{remark}{Remark} 
\newtheorem*{question}{Question}

\theoremstyle{definition}

\makeatletter
\@ifclassloaded{article}{
  \newtheorem{definicion}{Definición} [section]     % Se resetea en cada chapter
}{
  \newtheorem{definicion}{Definición} [chapter]     % Se resetea en cada chapter
}
\makeatother

\makeatletter
\@ifclassloaded{article}{
  \newtheorem{definition}{Definition} [section]     % Se resetea en cada chapter
}{
  \newtheorem{definition}{Definition} [chapter]     % Se resetea en cada chapter
}
\makeatother

\makeatletter
\@ifclassloaded{article}{
  \newtheorem{ex}{Example} [section]     % Se resetea en cada chapter
}{
  \newtheorem{ex}{Example} [chapter]     % Se resetea en cada chapter
}
\makeatother
\providecommand*{\exautorefname}{Example}
\providecommand{\definitionautorefname}{Definition}

\newtheorem*{note}{Note}

\newtheorem*{notacion}{Notación}
\newtheorem*{ejemplo}{Ejemplo}
\newtheorem*{ejercicio*}{Ejercicio}             % No numerado
\newtheorem{ejercicio}{Ejercicio} [section]     % Se resetea en cada section


% Modificar el formato de la numeración del teorema "ejercicio"
\renewcommand{\theejercicio}{%
  \ifnum\value{section}=0 % Si no se ha iniciado ninguna sección
    \arabic{ejercicio}% Solo mostrar el número de ejercicio
  \else
    \thesection.\arabic{ejercicio}% Mostrar número de sección y número de ejercicio
  \fi
}

\newtheorem{problem}{Problem} [section]     % Se resetea en cada section

\newtheorem*{solution}{Solution} 
%\newtheorem*{pro}{Proposition} 
\newtheorem*{claim}{Claim} 
\newtheorem*{conc}{Conclusion} 


% Modificar el formato de la numeración del teorema "ejercicio"
\renewcommand{\theproblem}{%
  \ifnum\value{section}=0 % Si no se ha iniciado ninguna sección
    \arabic{problem}% Solo mostrar el número de ejercicio
  \else
    \thesection.\arabic{problem}% Mostrar número de sección y número de ejercicio
  \fi
}


% \renewcommand\qedsymbol{$\blacksquare$}         % Cambiar símbolo QED

% Recuadro para las soluciones
%\newtcolorbox[auto counter, number within=section]{solbox}[2][]{colframe=black, colback=white, coltitle=black, boxrule=0.3mm, sharp corners, boxsep=1mm, left=2mm, right=2mm, top=2mm, bottom=2mm, title={Solution~\thetcbcounter}, #1}
\newtcolorbox[auto counter, number within=section]{solbox}[2][]{colframe=black, colback=white, coltitle=black, 
  boxrule=0.2mm, sharp corners, boxsep=1mm, left=2mm, right=2mm, top=2mm, bottom=2mm, halign=justify, frame hidden, breakable,
  title={}, #1}
%------------------------------------------------------------------------

% Paquetes para encabezados
\usepackage{fancyhdr}
\pagestyle{fancy}
\fancyhf{}

\newcommand{\helv}{ % Modificación tamaño de letra
\fontfamily{}\fontsize{12}{12}\selectfont}
\setlength{\headheight}{15pt} % Amplía el tamaño del índice


%\usepackage{lastpage}   % Referenciar última pag   \pageref{LastPage}
%\fancyfoot[C]{%
%  \begin{minipage}{\textwidth}
%    \centering
%    ~\\
%    \thepage\\
%    \href{https://losdeldgiim.github.io/}{\texttt{\footnotesize losdeldgiim.github.io}}
%  \end{minipage}
%}
\fancyfoot[C]{\thepage}
%\fancyfoot[R]{\href{https://losdeldgiim.github.io/}{\texttt{\footnotesize losdeldgiim.github.io}}}

% Define a subtitle command
%\newcommand{\subtitle}[1]{%
%  \posttitle{%
%    \newline\par\large#1\par\vskip0.5em
%  }
%}
%------------------------------------------------------------------------

% Conseguir que no ponga "Capítulo 1". Sino solo "1."
\makeatletter
\@ifclassloaded{book}{
  \renewcommand{\chaptermark}[1]{\markboth{\thechapter.\ #1}{}} % En el encabezado
    
  \renewcommand{\@makechapterhead}[1]{%
  \vspace*{50\p@}%
  {\parindent \z@ \raggedright \normalfont
    \ifnum \c@secnumdepth >\m@ne
      \huge\bfseries \thechapter.\hspace{1em}\ignorespaces
    \fi
    \interlinepenalty\@M
    \Huge \bfseries #1\par\nobreak
    \vskip 40\p@
  }}
}
\makeatother

%------------------------------------------------------------------------
% Paquetes de cógido
\usepackage{minted}
%\renewcommand\listingscaption{Código fuente}

\usepackage{fancyvrb}
% Personaliza el tamaño de los números de línea
\DefineVerbatimEnvironment{output}{Verbatim}{
  fontsize=\small,
  frame=none,
  xleftmargin=1em,
  baselinestretch=1,
  gobble=0
}
\renewcommand{\theFancyVerbLine}{\small\arabic{FancyVerbLine}}

% Estilo para C++
\newminted{cpp}{
    frame=lines,
    framesep=2mm,
    baselinestretch=1.2,
    linenos,
    escapeinside=||
}

% para minted
\definecolor{LightGray}{rgb}{0.95,0.95,0.92}
\setminted{
    linenos=true,
    stepnumber=5,
    numberfirstline=true,
    autogobble,
    breaklines=true,
    breakautoindent=true,
    breaksymbolleft=,
    breaksymbolright=,
    breaksymbolindentleft=0pt,
    breaksymbolindentright=0pt,
    breaksymbolsepleft=0pt,
    breaksymbolsepright=0pt,
    fontsize=\footnotesize,
    bgcolor=LightGray,
    numbersep=10pt
}


\usepackage{listings} % Para incluir código desde un archivo

\renewcommand\lstlistingname{Código Fuente}
\renewcommand\lstlistlistingname{Índice de Códigos Fuente}

% Definir colores
\definecolor{vscodepurple}{rgb}{0.5,0,0.5}
\definecolor{vscodeblue}{rgb}{0,0,0.8}
\definecolor{vscodegreen}{rgb}{0,0.5,0}
\definecolor{vscodegray}{rgb}{0.5,0.5,0.5}
\definecolor{vscodebackground}{rgb}{0.97,0.97,0.97}
\definecolor{vscodelightgray}{rgb}{0.9,0.9,0.9}

% Configuración para el estilo de C similar a VSCode
\lstdefinestyle{vscode_C}{
  backgroundcolor=\color{vscodebackground},
  commentstyle=\color{vscodegreen},
  keywordstyle=\color{vscodeblue},
  numberstyle=\tiny\color{vscodegray},
  stringstyle=\color{vscodepurple},
  basicstyle=\scriptsize\ttfamily,
  breakatwhitespace=false,
  breaklines=true,
  captionpos=b,
  keepspaces=true,
  numbers=left,
  numbersep=5pt,
  showspaces=false,
  showstringspaces=false,
  showtabs=false,
  tabsize=2,
  frame=tb,
  framerule=0pt,
  aboveskip=10pt,
  belowskip=10pt,
  xleftmargin=10pt,
  xrightmargin=10pt,
  framexleftmargin=10pt,
  framexrightmargin=10pt,
  framesep=0pt,
  rulecolor=\color{vscodelightgray},
  backgroundcolor=\color{vscodebackground},
}

%------------------------------------------------------------------------

% Comandos definidos
\newcommand{\bb}[1]{\mathbb{#1}}
\newcommand{\cc}[1]{\mathcal{#1}}
\newcommand{\bs}[1]{\boldsymbol{#1}}
\newcommand{\R}{\bb{R}}
\newcommand{\N}{\bb{N}}

% I prefer the slanted \leq
\let\oldleq\leq % save them in case they're every wanted
\let\oldgeq\geq
\renewcommand{\leq}{\leqslant}
\renewcommand{\geq}{\geqslant}

% Si y solo si
\newcommand{\sii}{\iff}

% MCD y MCM
\DeclareMathOperator{\mcd}{mcd}
\DeclareMathOperator{\mcm}{mcm}

% Signo
\DeclareMathOperator{\sgn}{sgn}
\DeclareMathOperator{\spp}{supp}
\DeclareMathOperator{\tr}{tr}

% Letras griegas
\newcommand{\eps}{\epsilon}
\newcommand{\veps}{\varepsilon}
\newcommand{\lm}{\lambda}

\newcommand{\ol}{\overline}
\newcommand{\ul}{\underline}
\newcommand{\wt}{\widetilde}
\newcommand{\wh}{\widehat}

\let\oldvec\vec
\renewcommand{\vec}{\overrightarrow}

% Derivadas parciales
\newcommand{\del}[2]{\frac{\partial #1}{\partial #2}}
\newcommand{\Del}[3]{\frac{\partial^{#1} #2}{\partial #3^{#1}}}
\newcommand{\deld}[2]{\dfrac{\partial #1}{\partial #2}}
\newcommand{\Deld}[3]{\dfrac{\partial^{#1} #2}{\partial #3^{#1}}}


\newcommand{\AstIg}{\stackrel{(\ast)}{=}}
\newcommand{\Hop}{\stackrel{L'H\hat{o}pital}{=}}

\newcommand{\red}[1]{{\color{red}#1}} % Para integrales, destacar los cambios.

% Método de integración
\newcommand{\MetInt}[2]{
    \left[\begin{array}{c}
        #1 \\ #2
    \end{array}\right]
}

% Declarar aplicaciones
% 1. Nombre aplicación
% 2. Dominio
% 3. Codominio
% 4. Variable
% 5. Imagen de la variable
\newcommand{\Func}[5]{
    \begin{equation*}
        \begin{array}{rrll}
            \displaystyle #1:& \displaystyle  #2 & \longrightarrow & \displaystyle  #3\\
               & \displaystyle  #4 & \longmapsto & \displaystyle  #5
        \end{array}
    \end{equation*}
}

% Probabilidad
\DeclareMathOperator{\argmax}{arg\,max}
\DeclareMathOperator{\argmin}{arg\,min}
%\newcommand{\P}{\mathsf{P}}
\newcommand{\E}{\mathsf{E}}
\newcommand{\Var}{\mathsf{Var}}
\newcommand{\Prob}{\mathsf{P}}

%------------------------------------------------------------------------

\usepackage{multicol}

\usepgflibrary{arrows.meta} % LaTeX and plain TeX and pure pgf
\usepgflibrary[arrows.meta] % ConTeXt and pure pgf
\usetikzlibrary{arrows.meta} % LaTeX and plain TeX when using TikZ
\usetikzlibrary[arrows.meta] % ConTeXt when using TikZ

\newcommand{\myparagraph}[1]{\paragraph{#1}\mbox{}\\}

\newcommand\digitstyle{\color{blue}}
\makeatletter
\newcommand{\ProcessDigit}[1]
{%
  \ifnum\lst@mode=\lst@Pmode\relax%
   {\digitstyle #1}%
  \else
    #1%
  \fi
}
\makeatother
\lstset{literate=
    {0}{{{\ProcessDigit{0}}}}1
    {1}{{{\ProcessDigit{1}}}}1
    {2}{{{\ProcessDigit{2}}}}1
    {3}{{{\ProcessDigit{3}}}}1
    {4}{{{\ProcessDigit{4}}}}1
    {5}{{{\ProcessDigit{5}}}}1
    {6}{{{\ProcessDigit{6}}}}1
    {7}{{{\ProcessDigit{7}}}}1
    {8}{{{\ProcessDigit{8}}}}1
    {9}{{{\ProcessDigit{9}}}}1
    {<=}{{\(\leq\)}}1,
    morestring=[b]",
    morestring=[b]',
    morecomment=[l]//,
}

\definecolor{codegreen}{rgb}{0,0.6,0}
\definecolor{codegray}{rgb}{0.5,0.5,0.5}
\definecolor{codepurple}{rgb}{0.58,0,0.82}
\definecolor{backcolour}{rgb}{0.95,0.95,0.92}

\definecolor{backcolor}{RGB}{30,30,30}
\definecolor{codeblue}{RGB}{153,255,255}
\definecolor{codepink}{RGB}{200,80,180}
\definecolor{normalblue}{RGB}{12,144,244}

\lstset{
    language=C++,
    morekeywords={int,char,ld,ll,double}
}

\lstdefinestyle{mystyle}{
    backgroundcolor=\color{backcolour},   
    commentstyle=\color{codegreen},
    keywordstyle=\color{magenta},
    numberstyle=\tiny\color{codegray},
    stringstyle=\color{codepurple},
    %identifierstyle=\color{blue},
    basicstyle=\ttfamily\footnotesize,
    breakatwhitespace=false,         
    breaklines=true,                 
    captionpos=b,                    
    keepspaces=true,                 
    numbers=left,                    
    numbersep=5pt,                  
    showspaces=false,                
    showstringspaces=false,
    showtabs=false,                  
    tabsize=2
}
%\usemintedstyle[C++]{default}
\lstset{style=mystyle}
